\section{怎样划分文章的段落}

我们已经知道，文章中的一段话往往是由几个意思相对独立而又从不同侧面集中表达一个中心意思的句子或句群所组成。人们习惯地把这种类型的几个句子或句群叫做一段话的几个层次，而把由几个层次所组成的表达一个中心意思的一段话叫做文章的段落，或者叫做文章的“自然段”。这种“自然段”，能清晰地反映文章的内在结构，使文章眉目清楚，节奏鲜明，便于读者阅读和理解，因为一个自然段只能有一个中心意思，而且各个自然段的意思一定要有内在联系。若干个意思集中而彼此间又有内在联系的“自然段”就构成了一篇完整的文章。因此，我们在写文章时，切切不可全文一逗到底，全无划分自然段的概念；或者表面上划分了自然段，而把一些互不相涉、杂七杂八的意思放在一段中；或者把一个完整的意思在这一段说一点，在那一段又说一点，显得前后杂乱无章；或者只注意了事物的外部联系，实际上段与段之间缺乏内部联系，逻辑混乱。所有这些毛病，都由于没有掌握好划分自然段的要领，应该努力改正。

在文章中，自然段具有“换行”的明显标志，这是不难弄清的。它之所以“换行”，在于表示文章思想内容在表达时由于转折、强调、间歇等情况所造成的文字的停顿。为了表达好中心思想，使文章条理清晰起见，作者把文章写成很多自然段。读者为了更全面、更概括地理解和掌握文章的思想内容，就把若干自然段归纳为相对完整的一个“意义段”。这种“意义段”，人们又把它叫做“结构段”、“大段”、“层次”（大层次）或“部分”等。“意义段”与“自然段”的关系十分密切。意义段着眼于思想内容的划分：自然段侧重于文字表达的需要。一般说来，几个自然段合起来描述一个完整的事件或者论述一个完整的意思，方可组成一个意义段；当然也有时候一个自然段思想内容正好反映了一个意义段的完整意思，那么，二者就完全相合了。这种情况，与一个自然段一般包含了几个意思相对独立的小层次（有时就是一个小层次）大致相同。因此，人们常说的划分段落，指的是划分意义段。

那么，我们应该怎样划分文章的段落（意义段）呢？

\textbf{1. 精读全文，是划分文章意义段的基础。}

我们知道，句子的组织，层次的确定，自然段的划分，叙述的先后等，都要服从于表现文章主题的需要。如果我们只是想当然地把若干自然段拼凑成一个意义段，这就象不看目标乱放箭一样，肯定不能合理地准确地划分好意义段。只有在精读全文的基础上，才能大致了解文章思想内容和理出文章的脉络，才有可能对文章的各个组成部分进行准确、深入的分析。所以，那种看一个自然段就动手划分意义段的做法是行不通的。

\textbf{2. 根据不同文体的特点来划分文章意义段。}

这一点，同划分自然段的层次大致相同。

一般说来，记叙文的分段，往往按照事件发展的时间顺序进行，如刘白羽的《长江三峡》(见于部编语文高中第一册)写天已微明到八点二十分之所见，八点五十分之所见，十点以后之所见，中午之所见，加上开头一段写头一天在长江上航行的印象，五个意义段条理分明；有的还可以按人物、事件的空间变换来分段，如鲁迅的《从百草园到三味书屋》(见于部编语文初中第一册)，有的还可以综合运用时间的推移和空间的变换来分段，如《为了六十一级阶级弟兄》(见于部编语文高中第一册)把同一时间不同地点发生的种种事情交织在文章之中，文章中的小标题就是划分意义段的标志；有的还可以按照作者认识发展的顺序划分，如杨朔的《荔枝蜜》(见于部编语文初中第二册)；有的还可以按照故事情节的发展阶段分，如巴尔扎克的《守财奴》(见于部编语文高中第三册)；有的还可以按照先总、后分、再总的方法分，如魏巍的《谁是最可爱的人》(见于部编语文初中第四册)，等等。

议论文的分段，原则上可按提出问题(提出论述的中心)、分析问题(用论据进行论证)、解决问题(结论)的程序来考虑。其结构方式大致有分总式、总分式、总、分、总式、并列式、推进式等。由于具体文章的章法千变万化，我们要根据具体文章作具体的分析。有的是提出问题，面面铺开的方式进行论述，如毛泽东的《我们的文艺是为什么人 的？》（见于部编语文高中第一册），有的是用逐层剖析的 方式来论证中心论点，如《反对自由主义》（见于部编语文 初中第三册）；有的是用逐一批驳的方式来批判论敌的谬 误，如鲁迅的《“丧家的”“资本家的乏走狗”》（见于部 编语文高中第四册）；有的是边分析边作结展开论证，如斯 大林的《悼列宁》（见于部编语文高中第一册）。一般说来， 文章中论述的一个方面，就是一个意义段。我们只要认真揣 摩作者的思路，理清文章的脉络，就能合理地划分意义段。 说明文的段落原则上应按说明的项目和程序来划分。

说明文主要是介绍客观事理的文章，其客观事理与议论文的客 观事理是大致相近的。不过，议论文重议论，重证明和辩 驳；说明文重介绍，道理讲得少。我们就从这种情况出发， 根据文章说明事物的形状、特点、制作或发生发展的过程、 用途、功能、危害以及事物的本质属性等方面的条理和相互 关系，去划分说明文的意义段，说明事物的特征和本质。

\textbf{3.按照文章原有的某些标志来划分意义段。}

有的文章按照事情的内容或事物发展的时间、顺序等等 加了小标题，如萧伯赞的《内蒙访古》（见于部编语文高中 第四册），全文分两部分，每部分都有小标题：“一段最古 的长城”；“在大青山下”。再如李四光的《人类的出现》 （见于部编语文高中第二册），除了开头的总说为第一意义 段外，下面四段都有明确的小标题，那么全文分五个意义段 就很明确了。

有的文章虽没有运用小标题，但在内容相对完整的部分 之间，用空行或用花点隔开，也是自然成段。例如鲁迅的《<呐喊>自序》（见于部编语文高中第四册）、峻青的《海滨仲夏夜》（见于部编语文初中第一册）等。有的在每个部分上标以“一、二、三、四……”等数字，就成为章节式的意义段，如鲁迅的《狂人日记》、峻青的《党员登记表》（均见高中第三册）。对于这类文章，一般可以不再另行划分段落。但是，如果各部分所表达的内容比较广泛全面，也可以把段落分得具体一些，如《党员登记表》；如果各部分所表达的内容比较窄，而不够大，也可以把段落分得更概括些，如《狂人日记》。

\textbf{4.文章意义的划分还可依据文中承上启下的过渡句或过渡段进行。}

过渡句或过渡段，表示文章内容的转移或转折，在文章中结构中起承上启下的连接作用，有助于划分文章意义段，有时可以成为我们划分段落的显著标志。例如《谁是最可爱的人》一文，第一部分包括文章开头的三个自然段，作者魏巍以抒情的笔调，明确指出志愿军战士是最可爱的人；第二部分包括“4——14”自然段，报道志愿军的英雄事迹，以具体生动的事例歌颂志愿军战士是最可爱的人，而划段的明显标志就是第四自然段：“让我还是来说一段故事吧。”这一句就是过渡句，它起了承上启下的作用。至于第十四自然段用“朋友们，用不着多举例，你已经可以了解我们的战士是怎样一种人，这种人有什么一种品质，他们的灵魂多么地美丽和宽广。······”对上面所写的感人事例作了总结，这就是过渡段，标志着上面这个意义段到此结束。最后一个自然段也就是全文的第三部分，用丰富的联想和强烈的激情，对全文画龙点睛，指出“他们确实是我们最可爱的人”。

应该注意，过渡句或过渡段在一般情况下应分在下面一段，如上述的第四自然段。但是如果过渡句或过渡段只起总结上文的作用，则应分在上段，如上述的第十四自然段。简言之，过渡句或过渡段只起启下的作用的，应分在下段；只总结上文的，应分在上段，既总结上文又引起下文的，则以分在下段为宜。

划分段落之后，还要概括出段意。自然段的段意是根据自然段的主要内容归纳出来的，意义段的段意是根据组成意义段的几个自然段的段意进一步归纳出来的。段意的语言要简要、明确，尽可能选用文章的原句。归纳段意可以采用以下几种方法：

1. 对于记叙人物、事物或事件的段落，我们应该先简炼地概括出人或事的具体情况，然后揭示出这些人或事的思
想意义。就是说，不仅要指出“记叙了什么”，而且还要指出“表现了什么”。例如鲁迅的《为了忘却的记念》（见于部编语文高中第三册）第四部分：写柔石等被捕遇害的经过，揭露反动派的罪行，表达对烈士的深切悼念。

2. 对于不具体记叙人物、事物或事件而超说明写作目的，介绍时代背景，交代事件发生的时间、地点、起因，承上启下等作用的段落，我们应从它写了什么或起了什么作用等方面来概括。例如《刑场上的婚礼》（见于部编语文初中第六册）的第一个意义段（也是全文的第一个自然段），其段意是：提出两个使读者惊讶的问题，引起下文。

3. 对于记叙文中议论和抒情的段落，我们可引原文直接写出它所表达的思想观点和感情，也可以根据原文进行归纳。例如《谁是最可爱的人》的第一个意义段（包括一至三自然段）的段意可引用原文表达。作者以奔放的激情歌颂志愿军战士“是我们最可爱的人”。又如《为了忘却的记念》第五部分的段意可以这样归纳说：作者抒发了对烈士的怀念，表达了对反动派必然灭亡、人民革命必然胜利的坚强信念。

4. 对于议论文的段意，我们一般可引原文中的中心论点或分论点来表达，有的则需要根据中心内容进行归纳。例如《反对党八股》（部编语文高中第四册）可分三个部分。第一部分（第一自然段）：指明分析党八股的“八大罪状”是为了“以毒攻毒”；第二部分（第二至第九自然段）：列举党八股的“八大罪状”，指出其危害、反动实质和克服的办法；第三部分（最后一个自然段）：总结全文，阐明反对党八股，树立马列主义新文风的重大意义。

5. 说明文的段意可以根据说明的内容来表达。例如《中国石拱桥》（部编语文初中第三册）可分四个部分。第一部分（一至三自然段）：总说石拱桥及其历史，特别说明在中国有悠久的历史；第二部分（第四、五自然段）：说明赵州桥的地理位置、历史情况和杰出的建筑构造；第三部分（第六、七、八自然段）：说明芦沟桥的地理位置、历史情况、建筑艺术和历史意义；第四部分（结尾两个自然段）：指出我国石拱桥有这样光辉成就的原因以及在社会主义制度下的光辉前景。

6. 各类文体的各个部分的首句、结句或段中主句，往往有提示、总结或突出中心的作用，归纳段意时要注意它们，有的还可以直接引用。例如茅盾的《风景谈》（部编语文高中第三册）的六个画面（每一画面构成一个意义段），每个画面之后的结句都与段意密切相关，我们在归纳段意时，当然不能忽略它们。
\newpage